\section{Experimental Design}
\label{sec:exp}

\paragraph{Target and methods.}
The target is the population mean $\mathbb E[Y]$. For every method we report a Wald interval at nominal 95\% level, $\hat\theta\pm1.96\sqrt{\widehat V}$, and compare to:
(i) \emph{classical}, the labeled-only mean;
(ii) \emph{PPI} with $\lambda=1$;
(iii) \emph{PPI++} with the plug-in optimal $\lambda$ (linear calibration on $f$; or on $(f,X)$ when covariates are used);
(iv) \emph{stratified PPI} (Fisch et al., 2024, $\lambda_k=1$), with quantile strata of $f$ and unlabeled-estimated stratum weights, using the linearized variance;
(v) \emph{cross-fitted calibration} (ours), with 5-fold cross-fitting and three learners: a cubic B-spline ridge model (primary), gradient boosting, and linear regression.
The primary learner was chosen before the main runs using development seeds disjoint from the reported ones. The rule was the lowest mean width ratio to PPI++ across the five synthetic scenarios, subject to nominal coverage. The spline calibrator won on average (0.967, versus 0.978 for boosting and 1.009 for linear), and boosting was marginally better only in the nonlinear scenario. We report all learners on the real data.

\paragraph{Real data.}
We use the nine pools distributed with \texttt{ppi-python} v0.2.3\footnote{\url{https://github.com/aangelopoulos/ppi_py}}, which contain labels and predictions for every unit. Seven pools give exact coverage checks: \texttt{galaxies}, \texttt{forest}, \texttt{gene\_expression}, \texttt{alphafold}, \texttt{census\_income}, \texttt{census\_healthcare}, and \texttt{census\_education}. Pools \texttt{plankton} and \texttt{ballots} have no labels on part of the unlabeled pool, so their target is undefined and we exclude them from coverage (a shift problem for \texttt{plankton}). Covariates are used only in the \emph{fX} regime; constant columns are dropped. Each pool is treated as the finite population. Labeled and unlabeled samples are drawn i.i.d.\ with replacement from it, so the sampling model is exactly the i.i.d.\ model of Section~\ref{sec:setting} and the target is the pool mean. Sample sizes are $n\in\{50,100,200,500,1000\}$ with $N/n\in\{1,4,16\}$. We use \REPSREAL{} replications per cell.

\paragraph{Synthetic data.}
Five data-generating processes have known $\mathbb E[Y]$ (Monte Carlo truth with $10^7$ draws). \emph{linear} has $m$ affine in $f$, so PPI++ should be efficient. \emph{nonlinear} has $m(f)=\sin(2f)+f/2$. \emph{binary} has a Bernoulli outcome and a miscalibrated probability predictor. \emph{xonly} has $m$ depending on a covariate not in $f$. \emph{heavy} is nonlinear with $t_3$ noise. We use \REPSSYN{} replications per cell over the same grid.

\paragraph{Theory checks.}
To test Propositions~\ref{prop:var} and Corollary~\ref{cor:when} without learner noise, we run oracle calibrators $\hat m=m+s\,h(f)$, where $h$ is a centred, unit-variance function and $s\in\{0,0.1,0.2,0.4\}$ gives $\mathrm{Var}(\hat m-m)=s^2$. Population constants come from one large draw.

\paragraph{Reproducibility.}
Seeds are fixed per cell and derived from the cell's indices, so every number in the paper is regenerated by \texttt{run\_main.sh}. Compute ran on a 4-CPU quota; the full pipeline took \TOTALTIME{}.
